% ============================================================================
\clearpage
\section*{v4.10.0-r51 Source-Only: A Medium under the Frame Map JB (Reading b)}
\addcontentsline{toc}{section}{v4.10.0-r51 Source-Only: A Medium under the Frame Map J\_B}
\label{sec:v4100-r51-jb-medium}
% ============================================================================

\paragraph{Purpose.}
The author asked whether the medium required by the OP5 comparison model could work for the photon frame and appear bent in the other
frames.  Reading (a), ordinary boosts, is recorded in v4.10.0-r50 (H120).  This layer tests reading (b): the bending is the frame map \(J_B\) of
convention F, the complex boost with rapidity \(i\pi/2\) along the M--T axis, by which T is the image of M.  It is an exploration with a method fixed
before the computation; it is not blind, since the sign change of light-like contributions under \(J_B\) was seen by hand.  \(J_B\) is kept apart from the
tachyon image of an action (rapidity \(i\pi\), H79) and from time reversal.  Earlier layers are not rewritten.

\noindent\textbf{Source-only status.}
This v4.10.0-r51 layer is source-only.  It does not build a PDF, does not assign a new DOI,
and does not authorize a public release.  The v5.3.0 Version DOI
\href{https://doi.org/10.5281/zenodo.24434403}{10.5281/zenodo.24434403} and the Concept DOI
\href{https://doi.org/10.5281/zenodo.19210728}{10.5281/zenodo.19210728} are unchanged by it.

\subsection*{v4.10.0-r51 results}
\label{subsec:v4100-r51-results}

\paragraph{Theoremlet v4.10.0-r51.1 (the fixed directions of \(J_B\)).}
With the axis \(e=x\), \(J_B\) maps \((t,x)\mapsto(-ix,-it)\) and fixes \(y,z\); it preserves \(\eta\).  Its eigendirections are the light-like \((1,1,0,0)\) with eigenvalue \(-i\), \((-1,1,0,0)\)
with \(+i\), and the transverse plane with eigenvalue 1.  A real timelike vector is mapped to \(i\) times a spacelike one.
\emph{Proof.} Exact evaluation (script \texttt{jb\_\allowbreak medium.py}).  \(\square\)

\paragraph{Theoremlet v4.10.0-r51.2 (invariant and sign-changing propagation tensors).}
A symmetric tensor \(S^{\mu\nu}\) is invariant under \(J_B\) exactly if its \((t,x)\) block is proportional to the metric, its transverse block arbitrary and the mixed blocks
zero; the speed along \(e\) is then exactly 1.  It changes sign exactly if its \((t,x)\) block has the form \(\bigl(\begin{smallmatrix}s&q\\q&s\end{smallmatrix}\bigr)\) and all other entries vanish; the light-like
stream terms \(l_\pm l_\pm\) along \(e\) are of this kind.
\emph{Proof.} Linear equations \(J_BSJ_B^{\mathsf T}=\pm S\) (script \texttt{jb\_\allowbreak medium.py}).  \(\square\)

\paragraph{Theoremlet v4.10.0-r51.3 (the cone limit in M and T).}
Two counter-streams along \(e\) with \(\kappa a=\kappa b=0.3\) give in M the effective metric \(\operatorname{diag}(1.6,-0.4,-1,-1)\), squared speeds between \(0.25\) and \(0.625\) and a cone limit of
\(0.803\); their \(J_B\) image is real, \(\operatorname{diag}(0.4,-1.6,-1,-1)\), with squared speeds between \(2.5\) and \(4\) and a value of \(4.66\).  A \(J_B\)-invariant medium slowed across \(e\) is luminal along \(e\)
and gives \(1.1995\) in both frames.  Under the boost with rapidity \(i\pi\) (H79) the counter-stream medium is unchanged and gives \(0.803\) in both frames.
\emph{Proof.} Roots of the dispersion relation over 82 directions and the cone limit of r47 over \(\alpha\in[0,1]\) (script \texttt{jb\_\allowbreak medium.py}).  \(\square\)

\paragraph{Remark v4.10.0-r51.4 (reading).}
Against the criteria fixed beforehand reading (b) does \emph{not carry}: no medium keeps the cone limit in M and in its \(J_B\) image at once.  The idea is
met literally but with the wrong sign: the structures of the photon frame are the fixed directions of \(J_B\), and a medium built from them bends from
narrowing in M into widening in T.  As a reading found afterwards, wider cones suit the tachyon frame T; and by H79 the dynamics in T is governed by
the tachyon image of the action (\(i\pi\)), under which the medium stays healthy, whereas \(J_B\) belongs to the readout of the register.  The bending by \(J_B\)
therefore concerns the description in T, not the propagation.  The vacuum bound of r50.3 holds in M independently of the reading.

\subsection*{v4.10.0-r51 frames}
\label{subsec:v4100-r51-frames}

M and T in the sense of convention F; \(J_B\) is the complex boost with rapidity \(i\pi/2\) along the M--T axis.  The tachyon image of an action (\(i\pi\), H79) is used only
as a control; time reversal is not used.

\subsection*{v4.10.0-r51 register entries}
\label{subsec:v4100-r51-register}

\begin{center}
\small
\begin{tabular}{@{}>{\raggedright\arraybackslash}p{0.34\textwidth}>{\raggedright\arraybackslash}p{0.58\textwidth}@{}}
\toprule
statement & status after v4.10.0-r51\\
\midrule
medium bent by \(J_B\) (reading b) & does not carry: narrowing in M becomes widening in T (r51.3)\\
fixed structures of \(J_B\) & light-like directions along the M--T axis and the transverse plane (r51.1)\\
OP5 & open\\
\bottomrule
\end{tabular}
\end{center}

\subsection*{v4.10.0-r51 notes}
\label{subsec:v4100-r51-notes}

\texttt{notes/\allowbreak research/\allowbreak op5\_\allowbreak jb\_\allowbreak medium\_\allowbreak 2026\_\allowbreak 11\_\allowbreak 23/} (method note commit 20c34729,
script \texttt{jb\_\allowbreak medium.py}, output in \texttt{run/}).

\subsection*{v4.10.0-r51 ledger and guard update}
\label{subsec:v4100-r51-ledger-guard}

\[
R_{\rm chk}^{v4100\text{-}r51}=(1,1,0)
\]
for (status recorded; computations checked in the cited script; a medium healthy in M and T under \(J_B\) not found).

The layer spends no guard:
\[
R_{\rm glob}^{(v4100\text{-}r51,+)}=(0,1,1,1,1),
\qquad
\Gamma_{\rm guard}^{v4100\text{-}r51}=(1,1,1,1,1),
\qquad
\beta_4=1/256\ \text{unspent}.
\]
The layer does not imply exact finite \(C5\) zero closure, exact all-readout terminality, \(C6/T2\), a global photon-clock/RH theorem,
or global all-readout closure.  OP1--OP5, C5 and C6/T2 remain open.
